\subsection{Completion of the endpoint-two alternative}

The endpoint-two surface admits a second noninteger root at large minima.
Two geometric bounds first confine the surface to a rectangle on which
positive polynomial identities give that root. A complete finite base
then handles the smaller fully distinct triples.

\begin{lemma}
\label{lem:endpoint-two-rectangle}
For ordered real exponents $4\leq a\leq b\leq c$, the equation
$h_C(2)=0$ requires $b<2a-2$. If $a\geq8$, it also requires
$c<9a/4-8$.
\end{lemma}

\begin{proof}
First put $b=2a-2+u$ and $c=b+v$, where $u,v\geq0$. Define
\begin{align*}
 A_0(u)&=2u^2+7(a-2)u+2(3a^2-14a+12),\\
 B_0(u)&=(a+2)u^3+(4a^2+10a-12)u^2\\
 &\quad +(4a^3+25a^2-48a+12)u+2(13a^3-28a^2+8a+8),\\
 D_0(u)&=2(a-2)(9a^3+24a^2-56a+16)\\
 &\quad +(33a^3+20a^2-208a+136)u
 +(20a^2+23a-62)u^2+4(a+2)u^3,\\
 E_0(u)&=2(3a^3-a^2-8a+4)+(5a^2+3a-10)u+(a+2)u^2.
\end{align*}
Direct expansion of the quotient determinant gives
\[
 h_C(2)=A_0B_0+\tfrac12(A_0B_0)'v+D_0v^2+E_0v^3,
\]
where the derivative is in $u$. Substitution $a=4+m$ makes every
coefficient of $A_0,B_0,D_0,E_0$ strictly positive for $m\geq0$.
In particular, $3a^2-14a+12=3m^2+10m+4$.
Thus $h_C(2)>0$ on $b\geq2a-2$, proving the first bound.

For $a\geq8$ and $a\leq b<2a-2$, parameterize the proposed maximum
tail by
\[
 a=8+m,\qquad b=\frac{a+(2a-2)t}{1+t},\qquad
 c=\frac{9a}{4}-8+v,\qquad m,t,v\geq0.
\]
The inverse $t=(b-a)/(2a-2-b)$ has positive denominator. The exact
identity is
\begin{align*}
 64(1+t)^3h_C(2)
 &=9m^6(1+2t)\bigl[14(1-t)^2+5t+t^2\bigr]\\
 &\quad+308m^5t(t-1)^2+P(m,t,v).
\end{align*}
The first two terms are nonnegative. The complete polynomial $P$,
listed in Appendix~\ref{sec:endpoint-two-identities}, has $84$ strictly
positive coefficients and constant $3\,014\,656$. Hence the right side
is positive, and so is $h_C(2)$. Together with the first bound, this
proves $c<9a/4-8$ on the zero surface.
\end{proof}

\begin{theorem}
\label{thm:endpoint-two-completion}
For ordered real $40\leq a\leq b\leq c$ with $h_C(2)=0$,
\[
 h_C(4)>0>h_C(5).
\]
Consequently the quotient has a noninteger root in $(4,5)$.
For every positive integer exponent triple with $h_C(2)=0$, the
quotient spectrum is nonintegral. The integer conclusion uses the
complete finite base specified below and the earlier small-minimum
certificates.
\end{theorem}

\begin{proof}
Write $s=a+b+c$, $p=ab+ac+bc$ and $r=abc$, and set
\[
 D_4=h_C(4)-3h_C(2),\qquad D_5=4h_C(2)-h_C(5).
\]
The full identities in Appendix~\ref{sec:endpoint-two-identities} give
the following positive expansions. For $a=8+m,b=8+t,c=8+u$,
$D_5$ has $44$ strictly positive coefficients and constant
$4\,629\,843$. Thus $D_5>0$ for arbitrary $a,b,c\geq8$.

Lemma~\ref{lem:endpoint-two-rectangle} confines an endpoint-two zero
at $a\geq40$ to $a\leq b<2a$ and $a\leq c<9a/4$. Put
\[
 a=40+m,\qquad b=\frac{a(1+2t)}{1+t},\qquad
 c=\frac{a(4+9u)}{4(1+u)},\qquad m,t,u\geq0.
\]
The inverse substitutions are $t=(b-a)/(2a-b)$ and
$u=4(c-a)/(9a-4c)$, with positive denominators. The polynomial
$64(1+t)^3(1+u)^3D_4$ has $112$ strictly positive coefficients and
constant $611\,305\,472$, all displayed in the appendix. Therefore
$D_4>0$. At $h_C(2)=0$ these two signs give $h_C(4)>0>h_C(5)$;
the intermediate value theorem supplies the root in $(4,5)$.

For positive integer triples, repeated exponents are nonintegral by
Theorem~\ref{thm:repeated}, and minima at most seven by
Corollary~\ref{cor:minimum-seven}. The only remaining base is
\[
 8\leq a\leq39,\qquad a<b<2a-2,\qquad b<c<\frac{9a}{4}-8.
\]
The largest permitted integer $c$ is $\lfloor(9a-33)/4\rfloor$.
This domain contains exactly $8\,658$ triples; its per-minimum counts
are listed in Appendix~\ref{sec:endpoint-two-identities}. Every endpoint
value is nonzero, with smallest absolute value $8\,208$.
The complete
\href{https://github.com/the-omega-institute/ideal-intersection-laplacian/blob/4afe151e1dc7193e6776cc1b08d3606aa804238e/results/endpoint-two-completion-base.csv}{integer certificate}
records $h_C(2)$ by integer Horner evaluation and
$\det(2I-C)=2h_C(2)$ by an independent fraction-free $6\times6$
Bareiss determinant for every triple. Thus no fully distinct endpoint-two
zero occurs in this necessary base. The real tail covers all larger
minima, completing the integer conclusion. The proof does not require
emptiness of the unbounded integer zero surface.
\end{proof}

\begin{corollary}
\label{cor:endpoint-two-parity-completion}
Every all-even exponent triple and every permutation of $(1,3,2)$ or
$(3,3,0)$ modulo four has a nonintegral ideal intersection graph,
without a size, gap or ratio bound.
\end{corollary}

\begin{proof}
Repeated triples and minima at most seven are already excluded. For the
others, Theorem~\ref{thm:uniform-low-root} gives a quotient root in
$(0,3)$. If all exponents are even, $C/2$ is an integer matrix, so
an integer spectrum of $C$ would have only even roots by the rational-root
theorem. In the two mixed classes, the shifted-root argument in
Theorem~\ref{thm:mixed-residue-classes} forces hypothetical odd integer
roots to be three modulo four. In either case the low root cannot be one,
so it must be two. Theorem~\ref{thm:endpoint-two-completion} contradicts
integrality. The complement and universal-class lifts transfer the result
to the ideal intersection graph. Exclusion of one in the mixed classes is
conditional on an integer spectrum.
\end{proof}
